本实验面向数字 0--9 的孤立词语音识别任务，
完成了从数据预处理、信号分析、特征提取、模型训练到噪声鲁棒性测试的完整流程。
实验首先对语音信号进行端点检测、预加重、分帧加窗和短时傅里叶分析，
并提取梅尔频率倒谱系数（MFCC）及其动态特征。
随后，分别采用 MFCC 统计特征结合 KNN、MFCC 序列特征结合 GMM，
以及基于 MFCC 动态特征图的 CNN 模型进行数字分类。
实验结果表明，在纯净语音条件下，GMM 方法取得了较高识别准确率；
在混合噪声条件下，识别性能随信噪比降低而下降。
通过对时域、频域、时频域特征和不同模型性能的对比分析，
实验验证了信号处理方法在语音识别任务中的有效性，
并讨论了小样本数据集下不同方法的优劣与局限。

